\chapter{Introduction}
\label{ch:intro}

\section{Objective}
Retail order-line data — one row per product sold within an order — sits at
the intersection of three classical machine-learning problem types:
descriptive analytics (what happened), predictive analytics (what will
happen), and diagnostic/prescriptive analytics (why it happened, and what
to do about it). This report uses the public \emph{Sample Superstore}
dataset to walk through all three, mirroring a realistic analytics
engagement:

\begin{enumerate}[label=\arabic*.]
  \item \textbf{Data quality and exploratory analysis} — establish that the
        data is trustworthy and characterise its statistical structure.
  \item \textbf{Feature engineering} — turn raw transactional fields into
        model-ready predictors, with explicit attention to information
        leakage.
  \item \textbf{Forecasting} — model and project the aggregate monthly
        sales time series.
  \item \textbf{Supervised learning} — model two line-level targets,
        \emph{profit} (regression) and \emph{is a line loss-making}
        (classification), and interpret what drives them.
\end{enumerate}

Each stage in this report is structured identically: a \textbf{Theory}
section derives or states the relevant statistical/ML machinery from first
principles, followed by a \textbf{Method} section describing how it was
applied to this dataset, then \textbf{Results} with the produced figures
and tables, and finally an \textbf{Interpretation} section connecting the
numbers back to a business reading. This mirrors, chapter for chapter, the
five numbered notebooks in \texttt{EY\_Training/ml\_superstore/notebooks/}
and the reusable modules in \texttt{src/}.

\section{Dataset}
The dataset contains $n = 9{,}994$ order lines spanning
\SI{4}{years} (2014--01--03 to 2017--12--30), $793$ unique customers and
$1{,}862$ unique products, organised into $5{,}009$ distinct orders across
the United States. Table~\ref{tab:intro-schema} summarises the raw schema.

\begin{table}[H]
  \centering
  \caption{Raw schema of the Sample Superstore CSV.}
  \label{tab:intro-schema}
  \begin{tabular}{@{}llp{7.5cm}@{}}
    \toprule
    \textbf{Field} & \textbf{Type} & \textbf{Description} \\
    \midrule
    Row ID, Order ID & id & Line and order identifiers \\
    Order Date, Ship Date & date & Order placement and shipment dates \\
    Ship Mode & categorical & Standard / Second / First Class, Same Day \\
    Customer ID, Customer Name & id/text & Buyer identity \\
    Segment & categorical & Consumer / Corporate / Home Office \\
    Country, City, State, Postal Code, Region & geography & U.S. location fields \\
    Product ID, Category, Sub-Category, Product Name & catalogue & 3-level product hierarchy \\
    Sales & numeric (\$) & Line revenue \\
    Quantity & numeric & Units sold \\
    Discount & numeric $[0,1)$ & Fractional discount applied \\
    Profit & numeric (\$) & Line profit (can be negative) \\
    \bottomrule
  \end{tabular}
\end{table}

\section{Software and reproducibility}
All analysis is implemented in Python 3.10 using \texttt{pandas},
\texttt{numpy}, \texttt{scikit-learn}, \texttt{statsmodels},
\texttt{matplotlib}/\texttt{seaborn}. Code lives in
\texttt{EY\_Training/ml\_superstore/src/} as importable modules
(\texttt{data\_loader}, \texttt{features}, \texttt{forecasting},
\texttt{modeling}, \texttt{pipeline}), exercised interactively by six
Jupyter notebooks and, for one-shot reproduction, by
\texttt{python -m src.pipeline}. Every figure reproduced in this report was
generated by that codebase and saved under \texttt{outputs/figures/}; every
number quoted in the text is read directly from the corresponding CSV in
\texttt{outputs/tables/} or the machine-written \texttt{outputs/REPORT.md}.
No figure or statistic in this document was computed by hand.

\section{Notation}
Throughout, $y_i$ denotes a target value for observation $i$,
$\hat y_i$ its model prediction, $\mathbf{x}_i \in \mathbb{R}^p$ its feature
vector, and $n$ the sample size of the set under discussion (which may be a
train, test, or full-data set, made explicit in context). For the monthly
sales series we write $y_t$, $t = 1, \dots, T$, for the value in month $t$.
$\mathbb{E}[\cdot]$, $\mathrm{Var}(\cdot)$ and $\mathrm{Cov}(\cdot,\cdot)$
denote expectation, variance and covariance under whatever probabilistic
model is being assumed at that point (population model for theory,
empirical/plug-in estimator when applied to data).
